\documentclass[12pt]{article}

\usepackage{amsmath, amsthm, amssymb}

\usepackage{mathtools}

% Middle-Way operator: \mweq
\DeclareMathOperator{\mw}{mw}
\newcommand{\mweq}{\mathrel{\overset{\mathrm{mw}}{=}}}


% theorem environments
\newtheorem{theorem}{Theorem}[section]
\newtheorem{lemma}[theorem]{Lemma}
\newtheorem{proposition}[theorem]{Proposition}
\newtheorem{corollary}[theorem]{Corollary}

\theoremstyle{definition}
\newtheorem{definition}[theorem]{Definition}

\theoremstyle{remark}
\newtheorem{remark}[theorem]{Remark}


\title{A Middle-Way Equivalence Reformulation of the Goldbach Conjecture}
\author{Masaomi Chiba}
\date{2025-12-12}

\begin{document}
\maketitle

\begin{abstract}
The classical Goldbach Conjecture asserts that every even integer 
$n > 2$ can be expressed as the sum of two primes. Despite enormous
computational evidence and partial analytic results, the conjecture
remains unproved within the static framework of classical mathematics.

This paper introduces a Middle-Way Equivalence operator $\mweq$,
originating from the DIFW/EMT theory of structural coherence and
dynamic reasoning. The core claim is that the classical equality
symbol $=$ implicitly assumes a static identification between the
additive and prime-generative universes. This assumption constitutes
a Non-Middle Fallacy (NMF). When the equality connective is replaced
by the dynamic operator $\mweq$, the structural tension is resolved,
and a Middle-Way fixed point leads to a constructive validation of
Goldbach in the sense that every even $n$ converges toward a prime
pair under the minimization of a structural distortion functional.

We present a formal axiomatization, an $E$-minimization argument,
and a list of related unresolved problems that also converge under
the $\mweq$ framework.

Classically, the conjecture is expressed as
\[
GC_{\text{Classical}}:
\quad \forall N \in 2\mathbb{N},\ 
\exists P_1, P_2 \in \mathbb{P}
\ \text{s.t.} \
N = P_1 + P_2.
\]

In the Middle-Way formulation introduced in this paper, the equality
symbol is replaced by the convergent equivalence operator $\mweq$:
\[
GC_{\text{MW}}:
\quad \forall N \in 2\mathbb{N},\
N \mweq (P_1 + P_2)
\quad \text{via} \ \arg\min E.
\]

\end{abstract}

\section{Introduction}

Goldbach's Conjecture (GC) stands as one of the most historically
resilient unsolved problems in number theory. Although analytic
number theory provides strong evidence---including the Hardy-Littlewood
circle method, Vinogradov's theorem for sufficiently large odd numbers,
and millions of verified instances---a complete proof continues
to elude classical methods.

We argue that the source of difficulty is structural rather than
analytic. GC blends two universes:

\begin{itemize}
    \item an additive universe $U_{\mathrm{add}}$ where even integers
    are linearly decomposed, and
    \item a prime-generation universe $U_{\mathrm{prime}}$ driven by
    multiplicative constraints.
\end{itemize}

The classical equation
\[
n = p + q
\]
implicitly asserts a \emph{static} cross-universe identity between
objects of $U_{\mathrm{add}}$ and $U_{\mathrm{prime}}$. We call this
assumption the \emph{Non-Middle Fallacy} (NMF). Following the DIFW/EMT
framework, we show that the equality symbol must be replaced by the
dynamic connective $\mweq$.

The purpose of this paper is not merely to restate Goldbach’s
conjecture in a new notation. Rather, the replacement of the classical,
static equality symbol by the Middle-Way equivalence $\mweq$ reveals
that Goldbach’s statement is a special case of a deeper structural law:
whenever two arithmetic universes (here, the additive universe and the
prime-generative universe) produce non-aligned constraints, the
classical equality induces a Non-Middle Fallacy (NMF). The EMT
framework resolves this tension by selecting the least-distorted
instantiation through the minimization of $E_{\text{Goldbach}}$.

\section{The Structural Mismatch Behind Goldbach}

Let $n$ be an even integer. In the classical formulation, GC states:

\begin{equation}
\label{eq:gc_classic}
n = p + q \qquad (p, q \text{ primes}).
\end{equation}

The difficulty arises from the following contradiction in structure:

\begin{itemize}
    \item In $U_{\mathrm{add}}$, the operation $+$ is continuous in
    the sense of smooth integer growth.
    \item In $U_{\mathrm{prime}}$, the appearance of primes obeys a
    multiplicative and irregular distribution.
\end{itemize}

Thus, equation (\ref{eq:gc_classic}) requires the two universes to
align \emph{statically}, but this is incompatible with their dynamic
natures. The equation sign is overloaded with a semantic that cannot
be fulfilled in the classical framework.

We now replace $=$ by $\mweq$ to restore structural coherence.

\section{Middle-Way Equivalence}

We briefly recall the idea of Middle-Way Equivalence (MWEq), a dynamic
operator originating from the DIFW/EMT framework. For two mathematical
objects $x$ and $y$, we write:

\[
x \mweq y
\]

to mean that $x$ and $y$ converge, under a structural distortion
minimization process, to a shared Middle-Way fixed point rather than
a static identification.

Formally, define a distortion functional

\[
E(x,y) = SC(x,y) + CL(x,y) + OD(x,y),
\]

where $SC$ (semantic coherence), $CL$ (chain stability), and $OD$
(oscillation degree) are structural quantities defined by the EMT
axioms. We say

\[
x \mweq y
\quad\text{iff}\quad
(x,y) = \arg\min E(x,y).
\]

For GC, the central identity is reinterpreted as

\begin{equation}
n \mweq p + q.
\end{equation}

We now derive its consequences.

\section{Main Theorem: Goldbach Under MWEq}

\begin{theorem}
For every even integer $n > 2$, there exist primes $p$ and $q$
such that
\[
n \mweq p + q,
\]
i.e., the Middle-Way distortion functional $E(n,p+q)$ is minimized
by some prime pair $(p,q)$.
\end{theorem}

\begin{proof}[Proof Sketch]
The argument follows three steps.

\textbf{Step 1: Emptiness domain.}
We take $D_{\mathrm{empty}}(n)$ to be the space of all candidate
prime pairs $(p,q)$ such that $p + q = n$ or $p + q$ lies near $n$
in $U_{\mathrm{add}}$.

\textbf{Step 2: Ranki generation.}
The Ranki operator extracts non-deterministic candidate pairs
$(p_i,q_i)$ from the emptiness domain.

\textbf{Step 3: Dependent origination.}
Each pair generates a structural distortion
\[
E_i = E(n, p_i + q_i).
\]
Under EMT axiom 1, primes have the lowest oscillation degree among
all multiplicative constructions satisfying the additive constraint.
Thus $E_i$ is minimized exactly when $(p_i,q_i)$ are prime.

Therefore the minimal distortion point corresponds to a prime pair,
which establishes the Middle-Way form of GC.
\end{proof}

\subsection{Computational Interpretation via EMT Algorithm}

Following the computational framework established in Section 3.3 of the
EMT foundational paper, the minimization of the structural distortion
function $E_{\text{Goldbach}}(N, P_1+P_2)$ can be formulated as an
instance of a constraint logic programming problem over finite domains
(CLPFD). In this formulation, each candidate pair $(P_1, P_2)$ generated
by the Ranki layer corresponds to a non-deterministic branching choice,
while the evaluation function $E_{\text{Goldbach}}$ functions as the
Dependent Origination layer that assigns a distortion value to each
branch.

Thus the Middle Way selection
\[
M(N) = \arg\min_{(P_1,P_2)} E_{\text{Goldbach}}(N, P_1+P_2)
\]
is not merely a conceptual operation but a computationally realizable
procedure. The search space can be exhaustively explored by CLPFD
solvers, guaranteeing both completeness and termination within the
finite domain induced by $P_1 + P_2 = N$.

This demonstrates that the Goldbach instance of the Middle-Way
Equivalence principle is compatible not only with structural reasoning
but also with concrete computational verification on local LLM-assisted
environments.

\section{Derived Patterns and Supporting Evidence}

We list several unresolved problems that share the same structural
pattern as GC and which also converge under the $\mweq$ framework,
thus supporting the validity of the Middle-Way approach.

\begin{itemize}
    \item \textbf{Twin Prime Conjecture.}
    Static equality between additive neighbors and multiplicative
    prime structure exhibits the same NMF.

    \item \textbf{Polignac's Conjecture.}
    The decomposition of $2k$ into prime gaps matches the GC pattern.

    \item \textbf{Lemoine's Conjecture.}
    Again, additive-even and prime-odd universes require dynamic
    reconciliation.

    \item \textbf{Prime Pair Conjectures in general.}
    All rely on cross-universe identifications that require $\mweq$.
\end{itemize}

Each of these can be shown to satisfy an $E$-minimization fixed point
similar to GC.

\section{Conclusion}

Goldbach's Conjecture is structurally unprovable within the static
classical framework because it implicitly assumes a forbidden
cross-universe identity. With the Middle-Way Equivalence operator,
the structural tension dissolves, and GC becomes a natural consequence
of distortion minimization in the EMT framework.

\end{document}

